\documentclass{article}
\usepackage{amsmath,amsthm,amssymb} % Required for math stuff
\usepackage{graphicx} % Required for inserting images


\newtheorem{theorem}{Theorem}[section]

\title{Week 8 Submission}
\author{Mengxiang Jiang}
\date{\today}



\begin{document}

\maketitle

\section{Introduction}
There's an often told story of Hippasus, a mathematician during 
Pythagoras's time, that found out that the length of the diagonal of 
a square was not commensurate with the length of its side. He was 
apparently drowned for this discovery by the Pythagoreans. At some 
level, I can sympathize with the Pythagoreans, since this proof
is deeply unsettling when I first read it. Like, rational numbers
can be made as small and accurate as you want, what do you mean
there's something it can't perfectly describe? And even after
learning real analysis, the way $\sqrt{2}$ is actually described
(whether it's a Dedekind cut, equivalence class of Cauchy
sequences, or something else) leaves much to be desired.
I much prefer the algebraic treatment of $\sqrt{2}$ instead,
where just like $i$, we just claim it's not a part of the rationals
and adjoin it, without making any claims about what it really is.
So my proof won't say that $\sqrt{2}$ is irrational, but rather
that the rational numbers do not have such a number.

\section{The Proof}
\begin{theorem}
$$\forall q \in \mathbb{Q}, q^2 \ne 2.$$
\end{theorem}

\begin{proof}
Assume for contradiction that there does exist a number $q$ in the rationals whose
square is 2. We use the property that all rational numbers can be written in the
most reduced form (proof of this property left to the reader) to write 
\[
    q = \frac{a}{b},
\] 
where $a, b \in \mathbb{Z}$, $b \ne 0$, and $a$ and $b$ are coprime.
Thus
\[
    q^2 = 2 = \frac{a^2}{b^2},
\]
which implies
\[
    2b^2 = a^2.
\]
We use the property that the product of even numbers is even and the product
of odd numbers is odd (proof left to the reader) to conclude that $a$ must be
even since $a^2$ is even. Thus we can write $a=2c$ for some $c \in \mathbb{Z}$.
We then have
\[
    q^2 = 2 = \frac{(2c)^2}{b^2},
\]
which implies
\[
    2b^2 = 4c^2.
\]
Dividing both sides by 2, we have
\[
    b^2 = 2c^2,
\]
which using the latter property again tells us $b$ is even. But this
contradicts the fact that $a$ and $b$ are coprime since they both share
a factor of 2.
\end{proof}



\end{document}
